\documentclass[11pt]{article}
\usepackage[a4paper,margin=25mm]{geometry}
\usepackage{amsmath,amssymb,amsthm,hyperref}
\newtheorem{theorem}{Theorem}
\newtheorem{lemma}{Lemma}
\title{Eventual admissible equal-sided sizes at factorial order}
\author{Proof with completed independent mathematical reviews}
\date{30 September 2026}
\begin{document}\maketitle
This note combines the short signed fair target, the polynomial-growth
Ore bound, and the new order-doubling positive-prefix completion.
It improves the earlier eventual-size threshold, without changing the
known asymptotic upper bound for the \emph{smallest} fair size.

\begin{theorem}\label{main}
There is an absolute constant $C$ such that, for every $n\ge2$, every
integer $M$ satisfying
\[
 \prod_{p\le n}p\mid M,
 \qquad M\ge\exp(Cn\log(n+1))
\]
is the common number of faces of a permutation-fair collection of
$n$ equiprobable finite dice with globally distinct face labels.
More precisely, there are two attainable common counts $A,A+P$ with
\[
 P=\prod_{p\le n}p,\quad P\mid A,
 \quad A+P\le\exp(Cn\log(n+1)).
\]
Every positive multiple of $P$ at least $A(A-P)/P$ is then attainable.
\end{theorem}

\begin{proof}
Write $T_a=\exp(aH)$ in the reduced integral-signature group, with
$H=X_1+\cdots+X_n$.
The reviewed constructive signed-signature theorem and its base-digit
compression, in \path{quantitative_eventual_sizes.tex}, give a signed
word for $T_P$ of length at most
\[
 L_*=2^{2n+3}n^4n!P=\exp(O(n\log(n+1))).
\]
This formula is the short-signed-target lemma of that note; the older
completion and conductor estimates there are not used here.

The reviewed growth/Ore lemma in \path{linear_logarithmic_ore_bound.tex}
says that a signed word of length $L$ is a right fraction $uv^{-1}$,
with both positive lengths at most
\[
 \left\lceil16D_nL\log(8D_nL)\right\rceil,
 \qquad D_n=\sum_{j=1}^n j\frac{n!}{(n-j)!}\le e n n!.
\]
Consequently $T_P=uv^{-1}$ with $|u|,|v|\le H_n$ for some
$H_n=\exp(O(n\log(n+1)))$.

Apply \path{order_doubling_prefix_completion.tex} to $v$.
There is a positive suffix $c$ with $vc=T_A$ and
\[
 A\le(H_n+1)\exp(O(n\log(n+1)))
    =\exp(O(n\log(n+1))).
\]
Here equations of words denote equations of signatures.
Then
$T_PT_A=uv^{-1}vc=uc$ is positive, so both $A$ and $A+P$ are
attainable common counts. The subset-factorial divisibility condition
forces $P\mid A$: restrict to any $p$ dice, for each prime $p\le n$,
and use $p!\mid A^p$.

Put $a=A/P$. Concatenating balanced fair words adds their common counts,
because $T_xT_y=T_{x+y}$.
Thus every nonnegative integer combination of $a$ and $a+1$, multiplied
by $P$, is attainable. Every integer $t\ge a(a-1)$ is such a combination:
write $t=qa+r$ with $0\le r<a$, so $q\ge a-1\ge r$ and
$t=(q-r)a+r(a+1)$.
Therefore every admissible count at least
$P a(a-1)=A(A-P)/P$ is attainable. Its bound is still
$\exp(O(n\log(n+1)))$, proving the theorem.
Finally a positive word is realized by assigning successive globally
distinct integers to its positions and interpreting equal letters as
faces of the corresponding equiprobable die.
\end{proof}

\paragraph{Scope and attribution.}
The theorem says that above a factorial-order threshold the elementary
prime divisibility conditions are sufficient for \emph{every} equal-sided
size. It does not distinguish $\exp(\Theta(n))$ from
$\exp(\Theta(n\log n))$ for the minimum total size, and does not prove
polynomial equalization of an arbitrary optimal family.
The order-doubling input uses the published Kuperberg--Lovett--Peled
$t$-wise permutation theorem. All other inputs are identified by their
separate proof notes; no novelty claim about the classical tools is made.
\end{document}
